\documentclass[11pt]{ctexart}

\usepackage{geometry}
\usepackage{amsmath,amssymb}   % 数学公式
\usepackage{graphicx}          % 插图
\usepackage{booktabs}          % 三线表 \toprule \midrule \bottomrule
\usepackage{multirow}          % 表格跨行单元格（如 SCF 列）
\usepackage{enumitem}          % 更好用的列表
\usepackage{xcolor}
\usepackage{fancyhdr}
\usepackage{titlesec}
\usepackage{indentfirst}        % ← 新增这行：让标题后第一段也缩进，不想让某一行锁进时，在这一段开头加\noindent
\usepackage[version=3]{mhchem}  % 化学式 \ce{}
\usepackage{float}              % 强制表格位置 [H]
\usepackage[hidelinks]{hyperref}

\geometry{a4paper,margin=1.5cm}
% \allowdisplaybreaks
\numberwithin{equation}{section}

\newcommand{\dd}{\mathrm{d}}
\newcommand{\pd}[2]{\frac{\partial #1}{\partial #2}}
\newcommand{\bra}[1]{\langle #1|}
\newcommand{\ket}[1]{|#1\rangle}
\newcommand{\eri}[1]{\!\left(#1\right)}

\title{spin flip up excitation TDA analytic gradient}
\author{王宏博}
\date{2026.6.14}

\begin{document}
\maketitle
\vspace{-2.2em}
\section{excited state analytic gradient basis}
激发态解析梯度可以写为基态能量的梯度加激发能的梯度，即
\begin{equation}
    \frac{d E_{exc}}{d\xi} = \frac{d E_{g}}{d \xi} + \frac{d \omega}{d \xi}
\end{equation}
基态能量的梯度在现有程序中容易计算，这里暂时不展开，以后补充这块内容。
要计算激发能的梯度，首先要得到激发能，对于 USF-TDA , Casida equation 为
\begin{equation}
    \mathbf{A}\mathbf{X} = \omega \mathbf{X}
\end{equation}
\begin{equation}
    \omega = \mathbf{X}^{T} \mathbf{A} \mathbf{X}
\end{equation}
\begin{equation}
    A_{ia,jb}^{\beta\alpha} 
    = \delta_{ij}F_{ab}^{\alpha} 
    - F_{ij}^{\beta}\delta_{ab} 
    - c_x (i^{\beta}j^{\beta}|b^{\alpha}a^{\alpha})
    + f_{pq,rs}^{xc,\sigma\tau}[\rho] 
\end{equation}
对于 collinear kernel 
\begin{equation}
    f_{pq,rs}^{xc,\sigma\tau}[\rho] = 0
\end{equation}
对于 multicollinear kernel 
\begin{equation}
    f_{pq,sr}^{xc,\sigma\tau}[\rho] 
    = \frac{\delta^2 E_{xc}}{\delta\rho_{pq}^{\sigma} \delta\rho_{sr}^{\tau}}
\end{equation}

对于 ROKS 参考态， Casida equation 与 USF-TDA 相同， $\mathbf{A}$ 同样与 USF-TDA 相同，注意 Fock 与 USF-TDA 不同，后面 Lagrange Multiplier method 中的约束也与 USF-TDA 不同。
\begin{equation}
    \mathbf{A}\mathbf{X} = \omega \mathbf{X}
\end{equation}
\begin{equation}
    \omega = \mathbf{X}^{T} \mathbf{A} \mathbf{X}
\end{equation}
\begin{equation}
    A_{ia,jb}^{\beta\alpha} 
    = \delta_{ij}F_{ab}^{\alpha} 
    - F_{ij}^{\beta}\delta_{ab} 
    - c_x (i^{\beta}j^{\beta}|b^{\alpha}a^{\alpha}) 
    + f_{pq,rs}^{xc,\sigma\tau}[\rho] 
\end{equation}
同样对于 collinear kernel
\begin{equation}
    f_{pq,rs}^{xc,\sigma\tau}[\rho] = 0
\end{equation}
对于 multicollinear kernel 
\begin{equation}
    f_{pq,sr}^{xc,\sigma\tau}[\rho] 
    = \frac{\delta^2 E_{xc}}{\delta\rho_{pq}^{\sigma} \delta\rho_{sr}^{\tau}}
\end{equation}


\section{UKS with USF-TDA analytic gradient}
这一节推导 UKS 参考态下 spin flip up TDA 激发能的解析梯度。所以后面的推导中所有 $i,j\in C$ ， $a,b\in V$。与 spin conversing U-TDDFT 的解析梯度类似，同样使用 Lagrangian Multiplier method ，同样的三个约束需要满足

\noindent
1. 正交 （由于考虑 TDA ，所以仅剩 $\mathbf{X}$）
\begin{equation}
    \mathbf{X}^{T}\mathbf{X} = 1
\end{equation}
2. Brillouin定理
\begin{equation}
    F_{ia}^{\sigma} = F_{ai}^{\sigma} = 0
\end{equation}
3. 正则轨道
\begin{equation}
    S_{pq}^{\sigma} = \delta_{pq}
\end{equation}
所以拉格朗日量可以写为
\begin{equation}
    \mathcal{L}[\mathbf{X}, \mathbf{C}, \Omega, \mathbf{Z}, \mathbf{W}]
    = \omega 
    - \Omega(\mathbf{X}^{T}\mathbf{X} - 1)
    + \sum_{ai\sigma}Z_{ai}^{\sigma}F_{ai}^{\sigma}
    - \sum_{pq\sigma}W_{pq}^{\sigma}(S_{pq}^{\sigma}-\delta_{pq})
\end{equation}
同样拉格朗日乘子法要求
\begin{equation}
    \frac{\partial \mathcal{L}}{\partial X_{ia}^{\sigma}} = 0,
    \quad \frac{\partial \mathcal{L}}{\partial C_{\mu p}^{\sigma}} = 0,
    \quad \frac{\partial \mathcal{L}}{\partial \Omega} = 0 ,
    \quad \frac{\partial \mathcal{L}}{\partial Z_{ai}^{\sigma}} = 0 ,
    \quad \frac{\partial \mathcal{L}}{\partial W_{pq}^{\sigma}} = 0 
\end{equation}
所以
\begin{equation}
    \frac{d \omega}{d \xi} 
    = \frac{\partial\mathcal{L}}{\partial \xi}
    + \frac{\partial\mathcal{L}}{\partial X_{ia}^{\sigma}}\frac{\partial X_{ia}^{\sigma}}{\partial \xi}
    + \frac{\partial\mathcal{L}}{\partial C_{\mu p}^{\sigma}}\frac{\partial C_{\mu p}^{\sigma}}{\partial \xi}
    + \frac{\partial\mathcal{L}}{\partial \Omega}\frac{\partial \Omega}{\partial \xi}
    + \frac{\partial\mathcal{L}}{\partial Z_{ai}^{\sigma}}\frac{\partial Z_{ai}^{\sigma}}{\partial \xi}
    + \frac{\partial\mathcal{L}}{\partial W_{pq}^{\sigma}}\frac{\partial W_{pq}^{\sigma}}{\partial \xi}
    = \frac{d \mathcal{L}}{d\xi}
\end{equation}
同样得到拉格朗日量对坐标的偏导数，由于仅考虑 TDA 所以其中一些项无需写为对称的形式。
% \begin{equation}
%     \begin{aligned}
%         \frac{\partial\mathcal{L}}{\partial\xi} 
%         &= \sum_{\mu\nu}\Bigg[
%             \sum_{\sigma}(T_{\mu\nu}^{\sigma}+Z_{\mu\nu}^{S,\sigma})\frac{\partial h_{\mu\nu}^{core}}{\partial\xi}
%             + \sum_{\sigma}(T_{\mu\nu}^{\sigma}+Z_{\mu\nu}^{S,\sigma})\frac{\partial g_{\mu\nu}^{\sigma}[\mathbf{D}]}{\partial\xi}
%             +\sum_{\sigma}(T_{\mu\nu}^{\sigma}+Z_{\mu\nu}^{S,\sigma})v_{\mu\nu}^{xc,\sigma,(\xi)}[\rho]\\
%             &\quad - c_x\sum_{\kappa\lambda}X_{\mu\nu}^{S,\beta\alpha}X_{\kappa\lambda}^{S,\beta\alpha}\frac{\partial (\mu\kappa|\lambda\nu)}{\partial\xi}
%             - c_x\sum_{\kappa\lambda}X_{\mu\nu}^{A,\beta\alpha}X_{\lambda\kappa}^{A,\beta\alpha}\frac{\partial(\mu\lambda|\kappa\nu)}{\partial\xi}
%             - \sum_{\sigma}\frac{\partial S_{\mu\nu}^{\sigma}}{\partial\xi} W_{\mu\nu}^{\sigma}
%         \Bigg]
%     \end{aligned}
% \end{equation}
\begin{equation}
    \begin{aligned}
        \frac{\partial\mathcal{L}}{\partial\xi} 
        &= \sum_{\mu\nu}\Bigg[
            \sum_{\sigma}(T_{\mu\nu}^{\sigma}+Z_{\mu\nu}^{S,\sigma})\frac{\partial h_{\mu\nu}^{core}}{\partial\xi}
            + \sum_{\sigma}(T_{\mu\nu}^{\sigma}+Z_{\mu\nu}^{S,\sigma})\frac{\partial g_{\mu\nu}^{\sigma}[\mathbf{D}]}{\partial\xi}
            +\sum_{\sigma}(T_{\mu\nu}^{\sigma}+Z_{\mu\nu}^{S,\sigma})v_{\mu\nu}^{xc,\sigma,(\xi)}[\rho]\\
            &\quad - c_x\sum_{\mu\nu}X_{\mu\nu}^{\beta\alpha}\frac{\partial K_{\mu\nu}^{\beta\alpha}[\mathbf{X}]}{\partial\xi}
            - \sum_{\sigma}\frac{\partial S_{\mu\nu}^{\sigma}}{\partial\xi} W_{\mu\nu}^{\sigma}
        \Bigg]
    \end{aligned}
\end{equation}
其中
\begin{equation}
    \begin{aligned}
        & T_{\mu\nu}^{\sigma} 
        = \sum_{ab}C_{\mu a}^{\alpha}T_{ab}C_{\nu b}^{\alpha}
        + \sum_{ij}C_{\mu i}^{\beta}T_{ij}C_{\nu}^{\beta}\\
        & T_{ab} = \sum_{i}X_{ia}X_{ib}\\
        & T_{ij} = -\sum_{a}X_{ia}X_{ja}\\
        & Z_{\mu\nu}^{S,\sigma} = \frac{1}{2}(Z_{\mu\nu}^{\sigma} + Z_{\nu\mu}^{\sigma})\\
        & g_{\mu\nu}^{\sigma}[\mathbf{V}] = \sum_{\kappa\lambda\tau}[(\mu\nu|\kappa\lambda)-c_x\delta_{\sigma\tau}(\mu\lambda|\kappa\nu)]V_{\kappa\lambda}\\
        & K_{\mu\nu}^{\beta\alpha}[\mathbf{X}] = \sum_{\kappa\lambda}(\mu\lambda|\kappa\nu)\sum_{jb}C_{\kappa j}^{\beta}C_{\lambda b}^{\alpha}X_{jb}\\
        & v_{\mu\nu}^{xc,\sigma,(\xi)}[\rho] = \frac{\partial v_{\mu\nu}^{xc,\sigma}[\rho]}{\partial \xi}\\
        & f_{\mu\nu,\kappa\lambda}^{xc,\sigma\tau,(\xi)}[\rho] = \frac{\partial f_{\mu\nu,\kappa\lambda}^{xc,\sigma\tau}[\rho]}{\partial\xi}
    \end{aligned}
\end{equation}
同样定义
\begin{equation}
    G_{\mu\nu}^{\sigma}[\mathbf{V}] = g_{\mu\nu}^{\sigma}[\mathbf{V}] + f_{\mu\nu}^{xc,\sigma}[\mathbf{V}]
\end{equation}
同样计算 $Q_{pq}^{\sigma}$
\begin{equation}
    \begin{aligned}
        & Q_{i' j'}^{\alpha'} = 2G_{i' j'}[\mathbf{T}]\\
        & Q_{a' b'}^{\alpha'} = 2\left(
            \sum_{a}T_{aa'}F_{ab'}^{\alpha'}
            - c_x\sum_{i}X_{ia'}K_{ib'}[\mathbf{X}]
        \right)\\
        & Q_{i' a'}^{\alpha'} = 2G_{i' a'}[\mathbf{T}]\\
        & Q_{a' i'}^{\alpha'} = 2\left(
            \sum_{a}T_{aa'}F_{a i'}^{\alpha'} 
            - c_x\sum_{i}X_{ia'}K_{ii'}[\mathbf{X}]
        \right)\\
        & Q_{i' j'}^{\beta'} = 2\left(
            \sum_{j}T_{i'j}F_{j'j}^{\beta'}
            + G_{i'j'}^{\beta'}[\mathbf{T}]
            - c_x \sum_{a}X_{i'a}K_{j'a}[\mathbf{X}]
        \right)\\
        & Q_{a' b'}^{\beta'} = 0\\
        & Q_{i' a'}^{\beta'} = 2\left(
            \sum_{j}T_{i'j}F_{a' j}^{\beta'} 
            + G_{i' a'}^{\beta'}[\mathbf{T}]
            - c_x\sum_{a}X_{i' a}K_{a'a}[\mathbf{X}]
        \right)\\
        & Q_{a' i'}^{\beta'} = 0
    \end{aligned}
\end{equation}
然后给出相同形式的 $\sum_{\mu'}\frac{\partial\mathcal{L}}{\partial C_{\mu' p'}^{\sigma'}}C_{\mu'q'}^{\sigma'}$ ，其中 $\sigma$ 取 $\alpha$ 和 $\beta$
\begin{equation}
    \begin{aligned}
        & \sum_{\mu'}\frac{\partial\mathcal{L}}{\partial C_{\mu' i'}^{\sigma'}}C_{\mu' j'}^{\sigma'}
        = Q_{i' j'}^{\sigma'}
        + 2G_{i' j'}^{\sigma'}[\mathbf{Z}^S]
        - 2W_{i' j'}^{\sigma'}\\
        & \sum_{\mu'}\frac{\partial\mathcal{L}}{\partial C_{\mu' a'}^{\sigma}}C_{\mu' b'}^{\sigma'}
        = Q_{a' b'}^{\sigma'} 
        - 2W_{a' b'}^{\sigma'}\\
        & \sum_{\mu'}\frac{\partial\mathcal{L}}{\partial C_{\mu' i'}^{\sigma'}}C_{\mu' a'}^{\sigma'}
        = Q_{i' a'}^{\sigma'}
        + \sum_{a}Z_{a i'}^{\sigma'}F_{a' a}^{\sigma'}
        + 2G_{i' a'}^{\sigma'}[\mathbf{Z}^S]
        - 2W_{i' a'}^{\sigma'}\\
        & \sum_{\mu'}\frac{\partial\mathcal{L}}{\partial C_{\mu' a'}^{\sigma'}}C_{\mu' i'}^{\sigma'}
        = Q_{a' i'}^{\sigma'}
        + \sum_{i}Z_{a' i}^{\sigma'}F_{i' i}^{\sigma'}
        - 2W_{a' i'}^{\sigma'}
    \end{aligned}
\end{equation}
相同形式的 Z vector equation
\begin{equation}
    - (Q_{i'a'}^{\sigma'} - Q_{a'i'}^{\sigma'})
    = \sum_{a}Z_{ai'}^{\sigma'}F_{a'a}^{\sigma'} 
    + 2G_{i'a'}^{\sigma'}[\mathbf{Z}^S]
    - \sum_{i}Z_{a'i}^{\sigma'}F_{i'i}^{\sigma'}
\end{equation}
相同形式的 $W_{pq}^{\sigma}$
\begin{equation}
    \begin{aligned}
        & W_{i'j'}^{\sigma'} = \frac{1}{2}Q_{i'j'}^{\sigma'} + G_{i'j'}^{\sigma'}\\
        & W_{a'b'}^{\sigma'} = \frac{1}{2}Q_{a'b'}^{\sigma'}\\
        & W_{i'a'}^{\sigma'} = \frac{1}{2}Q_{i'a'}^{\sigma'} + \frac{1}{2}\sum_{a}Z_{ai'}^{\sigma'}F_{a'a}^{\sigma'} + G_{i'a'}^{\sigma'}\\
        & W_{a'i'}^{\sigma'} = \frac{1}{2}Q_{a'i'}^{\sigma'} + \frac{1}{2}\sum_{i}Z_{a'i}^{\sigma'}F_{i'i}^{\sigma'}
    \end{aligned}
\end{equation}

\section{ROKS with USF-TDA analytic gradient}
结合前一节和 spin conversing excitation 中 ROKS 参考态下的推导，这一节的公式基本都能顺利得到，使用 Lagrangian Multiplier method ，同样有三个约束 

\noindent
1. 正交（由于考虑 TDA ，所以仅剩 $\mathbf{X}$ ）
\begin{equation}
    \mathbf{X}^{T}\mathbf{X} = 1
\end{equation}
2. Brillouin 定理发生变化
\begin{equation}
    F_{ai}^{\alpha}+F_{ai}^{\beta} = 0
    \quad F_{at}^{\alpha} = 0
    \quad F_{it}^{\beta} = 0
\end{equation}
3. 正则轨道
\begin{equation}
    S_{pq} = \delta_{pq}
\end{equation}
所以拉格朗日量可以写为
\begin{equation}
    \mathcal{L}[\mathbf{X},\mathbf{C},\Omega,\mathbf{Z},\mathbf{W}]
    = \omega 
    - \Omega(\sum_{ia}X_{ia}X_{ia}-1)
    + \sum_{ai}Z_{ai}(F_{ai}^{\alpha}+F_{ai}^{\beta})
    + \sum_{at}Z_{at}F_{at}^{\alpha}
    + \sum_{ti}Z_{ti}F_{ti}^{\beta}
    - \sum_{pq}W_{pq}(S_{pq}-\delta_{pq})
\end{equation}
同样拉格朗日乘子法要求
\begin{equation}
    \frac{\partial \mathcal{L}}{\partial X_{ia}} = 0,
    \quad \frac{\partial \mathcal{L}}{\partial C_{\mu p}} = 0,
    \quad \frac{\partial \mathcal{L}}{\partial \Omega} = 0 ,
    \quad \frac{\partial \mathcal{L}}{\partial Z_{ai}} = 0 ,
    \quad \frac{\partial \mathcal{L}}{\partial W_{pq}} = 0 
\end{equation}
所以
\begin{equation}
    \frac{d \omega}{d \xi} 
    = \frac{\partial\mathcal{L}}{\partial \xi}
    + \frac{\partial\mathcal{L}}{\partial X_{ia}^{\sigma}}\frac{\partial X_{ia}^{\sigma}}{\partial \xi}
    + \frac{\partial\mathcal{L}}{\partial C_{\mu p}^{\sigma}}\frac{\partial C_{\mu p}^{\sigma}}{\partial \xi}
    + \frac{\partial\mathcal{L}}{\partial \Omega}\frac{\partial \Omega}{\partial \xi}
    + \frac{\partial\mathcal{L}}{\partial Z_{ai}^{\sigma}}\frac{\partial Z_{ai}^{\sigma}}{\partial \xi}
    + \frac{\partial\mathcal{L}}{\partial W_{pq}^{\sigma}}\frac{\partial W_{pq}^{\sigma}}{\partial \xi}
    = \frac{d \mathcal{L}}{d\xi}
\end{equation}
同样得到拉格朗日量对坐标的偏导数，由于仅考虑 TDA 所以其中一些项无需写为对称的形式。
\begin{equation}\label{eq:ROKS USF partial L}
    \begin{aligned}
        \frac{\partial\mathcal{L}}{\partial\xi} 
        &= \sum_{\mu\nu}\Bigg[
            \sum_{\sigma}(T_{\mu\nu}^{\sigma}+Z_{\mu\nu}^{S})\frac{\partial h_{\mu\nu}^{core}}{\partial\xi}
            + \sum_{\sigma}(T_{\mu\nu}^{\sigma}+Z_{\mu\nu}^{S})\frac{\partial g_{\mu\nu}^{\sigma}[\mathbf{D}]}{\partial\xi}
            +\sum_{\sigma}(T_{\mu\nu}^{\sigma}+Z_{\mu\nu}^{S})v_{\mu\nu}^{xc,\sigma,(\xi)}[\rho]\\
            &\quad - c_x\sum_{\mu\nu}X_{\mu\nu}^{\beta\alpha}\frac{\partial K_{\mu\nu}^{\beta\alpha}[\mathbf{X}]}{\partial\xi}
            - \sum_{\sigma}\frac{\partial S_{\mu\nu}}{\partial\xi} W_{\mu\nu}
        \Bigg]
    \end{aligned}
\end{equation}
其中
\begin{equation}
    \begin{aligned}
        & T_{\mu\nu}^{\sigma} 
        = \sum_{ab}C_{\mu a}^{\alpha}T_{ab}C_{\nu b}^{\alpha}
        + \sum_{ij}C_{\mu i}^{\beta}T_{ij}C_{\nu}^{\beta}\\
        & T_{ab} = \sum_{i}X_{ia}X_{ib}\\
        & T_{ij} = -\sum_{a}X_{ia}X_{ja}\\
        & Z_{\mu\nu}^{S} 
        = \frac{1}{2}(Z_{\mu\nu} + Z_{\nu\mu})
        = Z_{\nu\mu}^{S}\\
        & Z_{\mu\nu} 
        = \sum_{ai}C_{\mu i}C_{\nu a}Z_{ai}
        + \sum_{at}C_{\mu t}C_{\nu a}Z_{ai}
        + \sum_{ti}C_{\mu t}C_{\nu i}Z_{ti}\\
        & g_{\mu\nu}^{\sigma}[\mathbf{V}] = \sum_{\kappa\lambda\tau}[(\mu\nu|\kappa\lambda)-c_x\delta_{\sigma\tau}(\mu\lambda|\kappa\nu)]V_{\kappa\lambda}\\
        & K_{\mu\nu}^{\beta\alpha}[\mathbf{X}] = \sum_{\kappa\lambda}(\mu\lambda|\kappa\nu)\sum_{jb}C_{\kappa j}^{\beta}C_{\lambda b}^{\alpha}X_{jb}\\
        & v_{\mu\nu}^{xc,\sigma,(\xi)}[\rho] = \frac{\partial v_{\mu\nu}^{xc,\sigma}[\rho]}{\partial \xi}\\
        & f_{\mu\nu,\kappa\lambda}^{xc,\sigma\tau,(\xi)}[\rho] = \frac{\partial f_{\mu\nu,\kappa\lambda}^{xc,\sigma\tau}[\rho]}{\partial\xi}
    \end{aligned}
\end{equation}
同样定义
\begin{equation}
    G_{\mu\nu}^{\sigma}[\mathbf{V}] = g_{\mu\nu}^{\sigma}[\mathbf{V}] + f_{\mu\nu}^{xc,\sigma}[\mathbf{V}]
\end{equation}
同样计算 $\sum_{\mu'}\frac{\partial\mathcal{L}}{\partial C_{\mu' p'}^{\sigma'}}C_{\mu'q'}^{\sigma'}=0$。首先计算 $Q_{pq}^{\sigma'}$
\begin{equation}
    \begin{aligned}
        & Q_{i' j'}^{\alpha'} = 2G_{i' j'}[\mathbf{T}]\\
        & Q_{i' t'}^{\alpha'} = 2G_{i't'}[\mathbf{T}]\\
        & Q_{i' a'}^{\alpha'} = 2G_{i' a'}[\mathbf{T}]\\
        & Q_{t' i'}^{\alpha'} = 2G_{t'i'}[\mathbf{T}]\\
        & Q_{t' u'}^{\alpha'} = 2G_{t'u'}[\mathbf{T}]\\
        & Q_{t' a'}^{\alpha'} = 2G_{t'a'}[\mathbf{T}]\\
        & Q_{a' i'}^{\alpha'} = 2\left(
            \sum_{a}T_{aa'}F_{a i'}^{\alpha'} 
            - c_x\sum_{i}X_{ia'}K_{ii'}[\mathbf{X}]
        \right)\\
        & Q_{a' t'}^{\alpha'} = 2\left(
            \sum_{a}T_{aa'}F_{at'}^{\alpha'}
            - c_x\sum_{i}X_{ia'}K_{it'}[\mathbf{X}]
        \right)\\
        & Q_{a' b'}^{\alpha'} = 2\left(
            \sum_{a}T_{aa'}F_{ab'}^{\alpha'}
            - c_x\sum_{i}X_{ia'}K_{ib'}[\mathbf{X}]
        \right)\\
        & Q_{i' j'}^{\beta'} = 2\left(
            \sum_{j}T_{i'j}F_{j'j}^{\beta'}
            + G_{i'j'}^{\beta'}[\mathbf{T}]
            - c_x \sum_{a}X_{i'a}K_{j'a}[\mathbf{X}]
        \right)\\
        & Q_{i' t'}^{\beta'} = 2\left(
            \sum_{j}T_{i'j}F_{t'j}^{\beta'}
            + G_{i't'}^{\beta'}[\mathbf{T}]
            - c_x\sum_{a}X_{i'a}K_{t'a}[\mathbf{X}]
        \right)\\
        & Q_{i' a'}^{\beta'} = 2\left(
            \sum_{j}T_{i'j}F_{a' j}^{\beta'} 
            + G_{i' a'}^{\beta'}[\mathbf{T}]
            - c_x\sum_{a}X_{i' a}K_{a'a}[\mathbf{X}]
        \right)\\
        & Q_{t'i'}^{\beta'} = 0\\
        & Q_{t'u'}^{\beta'} = 0\\
        & Q_{t'a'}^{\beta'} = 0\\
        & Q_{a' i'}^{\beta'} = 0\\
        & Q_{a't'}^{\beta'} = 0\\
        & Q_{a' b'}^{\beta'} = 0\\
    \end{aligned}
\end{equation}
然后得到
\begin{equation}
    \begin{aligned}
        & \sum_{\mu'}\frac{\partial\mathcal{L}}{\partial C_{\mu' i'}}C_{\mu' j'} 
        = Q_{i'j'} 
        + 2G_{i'j'}[\mathbf{Z}^{S}] 
        - 2W_{i'j'}\\
        & \sum_{\mu'}\frac{\partial\mathcal{L}}{\partial C_{\mu'a'}}C_{\mu'b'} 
        = Q_{a'b'}
        - 2W_{a'b'}\\
        & \sum_{\mu'}\frac{\partial\mathcal{L}}{\partial C_{\mu' t'}}C_{\mu' u'}
        = Q_{t'u'} 
        + 2G_{t'u'}^{\alpha'}[\mathbf{Z}^{S}] 
        - 2W_{t'u'} \\
        & \sum_{\mu'}\frac{\partial\mathcal{L}}{\partial C_{\mu'i'}}C_{\mu' a'}
        = Q_{i'a'} 
        + \sum_{a\sigma'}Z_{ai'}F_{aa'}^{\sigma'} 
        + \sum_{t}Z_{ti'}F_{ta'}^{\beta'}
        + 2G_{i'a'}[\mathbf{Z}^{S}] 
        - 2W_{i'a'}\\
        & \sum_{\mu'}\frac{\partial\mathcal{L}}{\partial C_{\mu' i'}}C_{\mu' t'}
        = Q_{i't'}
        + \sum_{a\sigma'}Z_{ai'}F_{at'}^{\sigma'}
        + \sum_{t}Z_{ti'}F_{tt'}^{\beta'}
        + 2G_{i't'}[\mathbf{Z}^{S}]
        - 2W_{i't'}\\
        & \sum_{\mu'}\frac{\partial\mathcal{L}}{\partial C_{\mu'a'}}C_{\mu' i'}
        = Q_{a'i'}
        + \sum_{i\sigma'}Z_{a'i}F_{i'i}^{\sigma'}
        + \sum_{t}Z_{a't}F_{i't}^{\alpha'}
        - 2W_{a'i'}\\
        & \sum_{\mu'}\frac{\partial\mathcal{L}}{\partial C_{\mu'a'}}C_{\mu' t'}
        = Q_{a't'}
        + \sum_{i\sigma'}Z_{a'i}F_{t'i}^{\sigma'}
        + \sum_{t}Z_{a't}F_{t't}^{\alpha'}
        - 2W_{a't'}\\
        & \sum_{\mu'}\frac{\partial\mathcal{L}}{\partial C_{\mu't'}}C_{\mu'i'}
        = Q_{t'i'}
        + \sum_{a}Z_{at'}F_{ai'}^{\alpha'}
        + \sum_{i}Z_{t' i}F_{i' i}^{\beta'}
        + 2G_{t'i'}^{\alpha'}[\mathbf{Z}^{S}]
        - 2W_{t'i'}\\
        & \sum_{\mu'}\frac{\partial\mathcal{L}}{\partial C_{\mu't'}}C_{\mu'a'}
        = Q_{t'a'}
        + \sum_{a}Z_{at'}F_{aa'}^{\alpha'}
        + \sum_{i}Z_{t' i}F_{a' i}^{\beta'}
        + 2G_{t'a'}^{\alpha'}[\mathbf{Z}^{S}]
        - 2W_{t'a'}
    \end{aligned}
\end{equation}
对于 ROKS 有 $W_{pq}=W_{qp}$ ，同样可以得到 Z vector equation
\begin{equation}
    \begin{aligned}
        & -(Q_{i'a'}-Q_{a'i'}) 
        = \sum_{i\sigma'}Z_{a'i}F_{i'i}^{\sigma'}
        + \sum_{t}Z_{a't}F_{i't}^{\alpha'}
        - \sum_{a\sigma'}Z_{ai'}F_{aa'}^{\sigma'} 
        - \sum_{t}Z_{ti'}F_{ta'}^{\beta'}
        - 2G_{i'a'}[\mathbf{Z}^{S}] \\
        & -(Q_{i't'}-Q_{t'i'}) 
        = \sum_{a}Z_{at'}F_{ai'}^{\alpha'}
        + \sum_{i}Z_{t' i}F_{i' i}^{\beta'}
        + 2G_{t'i'}^{\alpha'}[\mathbf{Z}^{S}]
        - \sum_{a\sigma'}Z_{ai'}F_{at'}^{\sigma'}
        - \sum_{t}Z_{ti'}F_{tt'}^{\beta'}
        - 2G_{i't'}[\mathbf{Z}^{S}]\\
        & -(Q_{t'a'}-Q_{a't'})
        = \sum_{i\sigma'}Z_{a'i}F_{t'i}^{\sigma'}
        + \sum_{t}Z_{a't}F_{t't}^{\alpha'}
        - \sum_{a}Z_{at'}F_{aa'}^{\alpha'}
        - \sum_{i}Z_{t' i}F_{a' i}^{\beta'}
        - 2G_{t'a'}^{\alpha'}[\mathbf{Z}^{S}]
    \end{aligned}
\end{equation}
求解 Z vector equation 可以得到 $Z_{ai}$, $Z_{at}$, $Z_{ti}$ ，进而可以得到 $W_{pq}$
\begin{equation}
    \begin{aligned}
        & W_{i'j'} = \frac{1}{2}Q_{i'j'} + G_{i'j'}[\mathbf{Z}^{S}]\\
        & W_{a'b'} = \frac{1}{2}Q_{a'b'}\\
        & W_{t'u'} = \frac{1}{2}Q_{t'u'} + G_{t'u'}^{\alpha'}[\mathbf{Z}^{S}]\\
        & W_{i'a'} = \frac{1}{2}Q_{i'a'} + \frac{1}{2}\sum_{a}Z_{ai'}F_{a'a}^{\sigma'} + \frac{1}{2}\sum_{t}Z_{ti'}F_{ta'}^{\beta'} + G_{i'a'}[\mathbf{Z}^{S}]\\
        & W_{i't'} = \frac{1}{2}Q_{i't'} + \frac{1}{2}\sum_{a\sigma'}Z_{ai'}F_{at'}^{\sigma'} + \frac{1}{2}\sum_{t}Z_{ti'}F_{tt'}^{\beta'} + G_{i't'}[\mathbf{Z}^{S}]\\
        & W_{a'i'} = \frac{1}{2}Q_{a'i'} + \frac{1}{2}\sum_{i}Z_{a'i}F_{i'i}^{\sigma'} + \frac{1}{2}\sum_{t}Z_{a't}F_{i't}^{\alpha'}\\
        & W_{a't'} = \frac{1}{2}Q_{a't'} + \frac{1}{2}\sum_{i\sigma'}Z_{a'i}F_{t'i}^{\sigma'} + \frac{1}{2}\sum_{t}Z_{a't}F_{t't}^{\alpha'}\\
        & W_{t'i'} = \frac{1}{2}Q_{t'i'} + \frac{1}{2}\sum_{a}Z_{at'}F_{ai'}^{\alpha'} + \frac{1}{2}\sum_{i}Z_{t' i}F_{i' i}^{\beta'} + G_{t'i'}^{\alpha'}[\mathbf{Z}^{S}]\\
        & W_{t'a'} = \frac{1}{2}Q_{t'a'} + \frac{1}{2}\sum_{a}Z_{at'}F_{aa'}^{\alpha'} + \frac{1}{2}\sum_{i}Z_{t' i}F_{a' i}^{\beta'} + G_{t'a'}^{\alpha'}[\mathbf{Z}^{S}]
    \end{aligned}
\end{equation}
有了 $W_{pq}$ 和 $Z_{pq}$ 带入 (\ref{eq:ROKS USF partial L}) ，然后加基态能量的梯度即可完成计算。



\end{document}
